\begin{table}[t]
\centering\footnotesize\setlength{\tabcolsep}{3pt}
\begin{tabular}{@{}lrrrrrl@{}}
\toprule
code & size & cand. & $t^*$ & colourings & $|F|$ at $t^*$ & by $k=1,\dots,5$ \\
\midrule
linear code \cite{baumert1971} & 343 & 0 & 0 & 0 & -- & 0, 0, 0, 0, 0 \\
Mathew--\"Osterg{\aa}rd \cite{mathew2017} & 350 & 7 & 7 & 64 & 1022--1118 & 7, 0, 0, 0, 0 \\
auxiliary-set search, $t=9$ & 363 & 20 & 18 & 96 & 2171--2308 & 2, 16, 2, 0, 0 \\
repaired auxiliary set, $t=10$ & 365 & 13 & 13 & 128 & 1582--1726 & 0, 13, 0, 0, 0 \\
repaired auxiliary set, $t=9$ & 366 & 12 & 12 & 64 & 1460--1561 & 0, 12, 0, 0, 0 \\
pipeline code 4 & 367 & 5 & 5 & 16 & 552--568 & 0, 5, 0, 0, 0 \\
$I_0$, our rebuild & 367 & 6 & 6 & 16 & 686--708 & 0, 6, 0, 0, 0 \\
pipeline code 0 & 367 & 6 & 6 & 16 & 686--708 & 0, 6, 0, 0, 0 \\
pipeline code 5 & 367 & 6 & 6 & 16 & 686--708 & 0, 6, 0, 0, 0 \\
pipeline code 6 & 367 & 6 & 6 & 16 & 688--708 & 0, 6, 0, 0, 0 \\
pipeline code 1 & 367 & 7 & 7 & 16 & 820--844 & 0, 7, 0, 0, 0 \\
pipeline code 2 & 367 & 7 & 7 & 16 & 822--848 & 0, 7, 0, 0, 0 \\
pipeline code 7 & 367 & 7 & 7 & 16 & 822--848 & 0, 7, 0, 0, 0 \\
pipeline, best & 367 & 8 & 8 & 16 & 956--984 & 0, 8, 0, 0, 0 \\
$I_0$, as printed \cite{polak2019} & 367 & 8 & 8 & 16 & 956--984 & 0, 8, 0, 0, 0 \\
pipeline code 3 & 367 & 8 & 8 & 16 & 956--984 & 0, 8, 0, 0, 0 \\
$X$ of \cite{gao2026}, as a code & 367 & 9 & 9 & 16 & 1098--1128 & 0, 9, 0, 0, 0 \\
ten-candidate code & 367 & 10 & 10 & 16 & 1240--1268 & 0, 10, 0, 0, 0 \\
$o=322$ auxiliary set & 367 & 10 & 10 & 16 & 1240--1268 & 0, 10, 0, 0, 0 \\
\bottomrule
\end{tabular}
\caption{The candidate count of every code we hold in $C_7^{\boxtimes 5}$. Candidates are the vertices with exactly one neighbour in the code; $t^*$ is the largest usable family (Corollary~\ref{cor:bip}); colourings are counted up to exchanging $\PH$ and $\PV$; $|F|$ ranges over them. Exhaustive for each file.}
\label{tab:codes}
\end{table}
